\section{Configuración de los experimentos}

Se compararán redes LSTM y transformadores con modelos de referencia lineales (OLS, Ridge, Lasso y Elastic Net) y no lineales (SVR con núcleo RBF, Random Forest y perceptrón multicapa). Dentro de cada escenario se utilizarán las mismas variables de entrada, ejemplos de entrenamiento e instantes de evaluación. El ajuste de hiperparámetros se realizará mediante validación temporal.

\begin{description}
  \item[Horizonte de una hora.] Se emplearán 24 horas de historia para estimar la generación de la hora siguiente.
  \item[Horizonte de 72 horas.] En la estrategia recursiva se incorporará cada pronóstico a la ventana de entrada para estimar el siguiente paso. En la estrategia directa, el modelo producirá las 72 estimaciones en una sola salida.
\end{description}

Se reportarán métricas globales y se analizará el error por paso del horizonte de 72 horas. El análisis considerará el solapamiento entre ventanas de evaluación y el costo computacional de los modelos.
